% Product 트랙 Day 6 슬라이드 본문 — 손으로 쓴 정본(한 프레임 = 한 쪽). Day5_슬라이드_body.tex 와 같은 형식.
% 래퍼: Day6_슬라이드.tex(슬라이드판) / Day6_슬라이드_노트.tex(노트판, 우측에 \note).
% 화면에는 결론만, 설명·근거·예외는 각 프레임의 \note{}에. 그림은 ../그림/Product_Day6/*.pdf.
% 모든 숫자는 Day6 워크북 완성 예시본(정본)과 day6_seed.py 에서. 시점 D+120~300 (2026-06-30 ~ 12-27).

\newcolumntype{L}[1]{>{\raggedright\arraybackslash}p{#1}}
\newcommand{\ra}{$\rightarrow$}
\newcommand{\til}{\textasciitilde}
\newcommand{\keymsg}[1]{\par\vfill{\color{mDarkTeal}\bfseries #1}\par}
\newcommand{\figcap}[1]{\par\smallskip{\footnotesize\color{black!60}#1}\par}
\newcommand{\fig}[2][0.66]{\centering\includegraphics[height=#1\textheight,width=\textwidth,keepaspectratio]{#2}\par}
\newcommand{\subhead}[1]{{\color{mDarkTeal}\bfseries #1}\par}
\newenvironment{tbl}[2][\small]{\begin{center}\usebeamercolor[fg]{normal text}\setlength{\tabcolsep}{4pt}\renewcommand{\arraystretch}{1.15}#1\begin{tabular}{#2}\toprule}{\bottomrule\end{tabular}\end{center}}

% =====================================================================
% 도입
% =====================================================================
\begin{frame}{Product 트랙 Day 6 — 시드: 무엇을 안 만들 것인가, 얼마를 받을 것인가}
\begin{itemize}
\item 8일 과정의 여섯째 날 · Product 트랙 2일차 · 시점 \textbf{D+120 \til{} D+300} (2026-06-30 \til{} 12-27)
\item 회사 상태: 5명 → \textbf{12명}, 시드 RCPS \textbf{12억}(D+180), 첫 유료 \textbf{3개사}, 현금 15.4억 / 22개월
\item 오늘의 문서 5종 ㉖\til{}㉚ — 그리고 세 사건: 조항 부재(D+186) · 퇴사(D+230) · 청구서 780만(D+280)
\end{itemize}
\keymsg{오늘의 네 질문 — 무엇을 만들지 않을 것인가 / AI 기능은 언제 된 것인가 / 무엇에 얼마를 받을 것인가 / 누구 돈으로 어떤 조건으로}
\note{표지. 어제 D+90에 쓴 다섯 문서(㉑\til{}㉕)가 오늘의 입력이다 — 학생들에게 ㉓ 항목 4·7, ㉕ 항목 6, ㉑ 항목 3을 손 닿는 곳에 두게 한다. 오늘은 0교시가 없다 — 1교시 90분, 마무리 30분. 사전 읽기는 정본 §3.2·§4.2\til{}4.4·§6.1\til{}6.3·§7.3 v1 열·§7.5(약 20분). 예상 소요 2분.}
\end{frame}

\begin{frame}[shrink=6]{오늘의 흐름}
\begin{columns}[T]
\begin{column}{0.46\textwidth}
\subhead{오전 — 결정과 범위, 그리고 AI의 완료 (제1\til{}3장)}
\begin{itemize}
\item \textbf{1교시} 90분 · 시드의 다섯 결정 · 세 시계 · v1.0/v1.1 · Appetite · 41 → 24 → 17 · Won't — 제1장 + 2.1\til{}2.4절
\item \textbf{2교시} 90분 · 게이트 · PR/FAQ · 사전등록 · AI 완료 정의 · 분자와 분모 · 데이터셋 — 2.5\til{}2.9절 + 3.1\til{}3.3절
\item \textbf{3교시} 90분 · 합격 기준·개정 · 오류 예산 · HITL · 밸류메트릭 · UE 두 열 — 3.4\til{}3.7절 + 4.1\til{}4.2절
\end{itemize}
\end{column}
\begin{column}{0.52\textwidth}
\subhead{오후 — 실습과 가격·자본이 번갈아}
\begin{itemize}
\item \textbf{4교시} 90분 · \textbf{실습 ①} ㉖ 범위 결정서 + ㉗ Eval·오류 예산·HITL — 워크북 §1\til{}2
\item \textbf{5교시} 75분 · 청구서 780만 · ARR 4,860 vs 6,060 · 계약 다섯 칸 · 텀시트·캡테이블·참가적·옵션풀 — 4.3\til{}4.6절 + 제5장
\item \textbf{6교시} 75분 · \textbf{실습 ②} ㉘ 프라이싱·UE · ㉙ PR/FAQ·사전등록 · ㉚ 텀시트·캡테이블 — 워크북 §3\til{}5
\item \textbf{마무리} 30분 · 학술 배경 · 읽을 자료 · Day7 예고
\end{itemize}
\end{column}
\end{columns}
\keymsg{오늘의 문서는 두 번 쓰인다 — v1.0을 쓰고, 사건이 나고, v1.1을 쓴다. 그 차이가 이 날의 학습이다}
\note{PM 트랙과 같은 6교시 틀(90·90·90·90·75·75 + 마무리 30 = 540). 실습 ①은 ㉖ 40분 + ㉗ 45분 + 안내 5분, 실습 ②는 ㉘ 30 + ㉙ 20 + ㉚ 25. 워크북 §0.1 "따라가는 순서"를 4교시 첫머리에 다시 읽는다. 오늘은 카드에 찾을 것이 많다 — 백로그 29건, 데이터셋 분자·분모, 파일럿 3사 제원, 청구서 내역, 텀시트 여덟 조항. 예상 소요 3분.}
\end{frame}

\begin{frame}{학습 목표 — 오늘이 끝나면}
\begin{columns}[T]
\begin{column}{0.5\textwidth}
\subhead{오전}
\begin{enumerate}
\item Appetite(기간)와 가용 공수(24)와 견적(41)을 \textbf{세 칸에 따로} 적고, Must 상한으로 정의를 줄여 자른다
\item 가용이 줄었을 때 \textbf{손실을 분해}하고(7.0 = 3.9 + 3.1) 남은 공수로 다시 나눈다
\item AI 기능의 완료를 \textbf{분자·분모·모집단·여집합}으로 쓰고, 0.1\%p에 기준을 고치지 않는다
\end{enumerate}
\end{column}
\begin{column}{0.48\textwidth}
\subhead{오후}
\begin{enumerate}\setcounter{enumi}{3}
\item 오류 예산의 \textbf{측정 원천이 놓치는 것}과 HITL의 \textbf{문구 둘}을 설명한다
\item 밸류메트릭을 4안 × 5기준으로 고르고, 청구서 780만 앞에서 \textbf{가설 넷과 계측}을 적는다
\item 캡테이블을 주식수로, 참가적 1x의 값(100억에서 9.6억)을 손으로, 없는 조항을 행으로 적는다
\end{enumerate}
\end{column}
\end{columns}
\keymsg{시드의 결정은 되돌리는 비용이 크다 — 6개월, 사고, 청구서, 다음 라운드. 그래서 숫자와 서명이 많다}
\note{여섯 목표 중 앞의 셋은 오전, 뒤의 셋은 오후. Day5의 keymsg "틀렸다는 것을 더 빨리 알기 위한 도구"의 시드판이 이 keymsg — 교재 1.1절. 예상 소요 2분.}
\end{frame}

% =====================================================================
\section{1교시 (90분) — 시드 단계의 제품 결정 · Appetite · 절단 · Won't}
% =====================================================================
\begin{frame}{이 교시에서 배우는 것}
\begin{itemize}
\item 시드의 다섯 결정 — 범위·완료 정의·가격·검증·자본, 그리고 되돌리는 비용 (1.1절)
\item 세 시계 — 런웨이·클로징·조건부 자금. D+149의 1,505만 (1.2절)
\item 문서가 두 번 쓰이는 이유 — v1.0, 사건, v1.1 (1.3절)
\item Appetite는 기간이다 · 견적과 절단 41 → 24 · Won't의 재검토 조건 · 두 번째 절단 24 → 17 (2.1\til{}2.4절)
\end{itemize}
\keymsg{돈이 들어오기 전에 범위를 자른다 — 텀시트는 "무엇을 만들 것인가"를 묻고, 41이면 24를 묻는다}
\note{교재 제1장 + 2.1\til{}2.4절. 90분 — 1.1절 12분(그림 1-1), 1.2절 10분(그림 1-2), 1.3절 8분, 2.1절 12분(그림 2-1), 2.2절 18분(그림 2-2), 2.3절 10분, 2.4절 15분(그림 2-3) + 정리 5분. 예상 소요 2분.}
\end{frame}

\begin{frame}{시드의 다섯 결정 — 프리시드에서는 가설, 시드에서는 숫자와 서명}
\begin{tbl}[\footnotesize]{L{16mm}L{36mm}L{26mm}L{40mm}L{9mm}}
결정 & 질문 & 프리시드 & 시드 & 문서 \\ \midrule
\textbf{범위} & 6개월 동안 무엇을 만들고 무엇을 만들지 않는가 & 트리의 잎 5 & 공수가 붙고 가용과 비교되고 \textbf{잘린다} & ㉖ \\
\textbf{완료 정의} & AI 판정 기능은 언제 "된 것"인가 & X-02 88.7\% & 데이터셋·합격 기준·오류 예산·사람의 자리 & ㉗ \\
\textbf{가격} & 무엇에, 얼마를 & 커널 한 줄 & 4안 비교 · 두 열 · 청구서 · 계약 다섯 칸 & ㉘ \\
\textbf{검증} & 틀렸다는 것을 어떻게 아는가 & 실험 카드 3 & 출시일의 글 · 미확정 FAQ · 사전등록 4 & ㉙ \\
\textbf{자본} & 누구 돈으로, 어떤 조건으로 & 런웨이 7 → 2.7 & 텀시트 여덟 조항 · 캡테이블 · 매각가 3구간 & ㉚ \\
\end{tbl}
\keymsg{공통점 — 되돌리는 비용이 크다: 6개월 · 사고 · 청구서 · 다음 라운드}
\note{교재 1.1절 표. Day5 1.1절의 네 결정(발견·범위·가격·지표) 중 오늘은 범위·가격이고, AI 제품에만 있는 완료 정의와 PM이 대개 못 보는 자본이 붙는다. 발견의 오류는 인터뷰 다섯 건이면 고쳐지지만 범위의 오류는 6개월이다. 예상 소요 6분.}
\end{frame}

\begin{frame}{시드 6개월의 시간표 — 문서 다섯, 게이트 셋, 사건 셋}
\fig[0.6]{fig_1_1_seed_timeline.pdf}
\keymsg{돈이 들어오기(D+180) 전에 범위를 잘랐다(D+130) — 순서가 옳다}
\note{그림 1-1. 위 — v1.0 다섯(D+130·137·151·158·172)과 게이트(D+200·240·300), 아래 — 사건(선투자·클로징·조항 부재·퇴사·청구서). 붉은 화살표 셋이 v1.1을 만든다 — 오늘 학생들이 세 번 v1.1을 쓴다. 정본 §4.2\til{}4.4, 워크북 §0.3. 예상 소요 6분.}
\end{frame}

\begin{frame}{세 시계 — 런웨이 · 클로징 · 조건부 자금}
\fig[0.6]{fig_1_2_cash_path.pdf}
\keymsg{D+149 잔액 1,505만 — 선투자가 15일 늦었으면 D+156 급여가 없었다. 15일이 회사의 수명이었다}
\note{그림 1-2 · 교재 1.2절 "작은 숫자로 먼저". 5,400만 ÷ 3,230만 = 1.67개월 → D+170에 0. 선투자 2억 D+150(전환, 12억에 포함), 잔여 10억 D+179, TIPS 1차년도 5억 D+217 [추가 설정]. 끝 15.4억/22개월(G-15). 텀시트 서명일(D+172)과 클로징(D+179) 사이에 급여일이 있다는 것, TIPS 조건(유료 5개사)이 범위·가격을 누른다는 것 — 두 가지를 강조. 예상 소요 8분.}
\end{frame}

\begin{frame}{문서가 두 번 쓰이는 이유 — v1.0이 있어야 v1.1이 측정이 된다}
\begin{tbl}[\footnotesize]{L{26mm}L{10mm}L{56mm}L{34mm}}
사건 & 문서 & v1.0이 있어서 알 수 있는 것 & v1.0이 없었다면 \\ \midrule
D+186 옵션풀 조항 부재 & ㉚ & 협상 1번이 있었고 구두 합의됐고 문서에 없다 — \textbf{문서화 실패} & "시드에서는 원래 안 정한다" \\
D+230 조현우 퇴사 & ㉖ & 손실 7.0 중 퇴사 2.2, 4.8은 그 전에 새어 나갔다 — 태산정밀 커스텀 & "개발자가 나가서 7 줄었다" \\
D+280 청구서 780만 & ㉘ & 원가 정렬 M을 \textbf{알고} 골랐고 계측을 안 넣었다 — 가설 4 판별 불가 & "재촬영 때문" — 정본을 읽은 답 \\
\end{tbl}
\keymsg{변경요청서(⑪)를 자기 자신에게 낸다 — 기준선이 있어야 변경이 있다}
\note{교재 1.3절. Day3 변경 통제의 자기 적용. 세 사건은 세 가지 다른 오류(문서화 실패·계획 외 착수·계측 미포함) — 하나로 뭉개면 원인이 사라진다. 사건 하나에 문서 전부를 고치지 않는다. 예상 소요 6분.}
\end{frame}

\begin{frame}{Appetite — 기간을 고정한다. 24도 41도 appetite가 아니다}
\fig[0.58]{fig_2_1_appetite.pdf}
\keymsg{Appetite = 6개월 · 가용 공수 = 24 mm · 견적 = 41 mm — 세 숫자를 세 칸에 따로 적는다}
\note{그림 2-1 · 교재 2.1절. Singer, Shape Up의 어휘를 정확히: 기간이 고정이면 사건(퇴사)에 줄어드는 것은 범위다. 흔한 오류 — Appetite 칸에 기능 목록, 6주 사이클 그대로, 회로 차단기 없이 "거의 다 됐으니 2주 연장", 예비 0. CTO를 100\%로 넣은 것이 v1.1에서 문제가 된다는 것을 예고. 예상 소요 8분.}
\end{frame}

\begin{frame}{견적과 절단 — 41 → 24 → 17}
\fig[0.58]{fig_2_2_backlog_cut.pdf}
\keymsg{견적은 개발자가 내고 PM은 깎지 않는다 — 깎는 것은 범위이고, 절단의 기술은 항목 빼기가 아니라 정의 줄이기다}
\note{그림 2-2 · 교재 2.2절. 29건 41.0 mm [추가 설정]. 상위 노드가 트리의 잎이 아닌 6건(고객 요청 4·기반 1·㉗ 1)이 절단의 첫 질문. 41 $-$ 24 = 17이 아니라 41 $-$ 18 = 23이 백로그 밖 — 예비 6은 항목이 아니다. 오른쪽 막대의 60\% 선이 "깨지는" 이유는 다음다음 프레임. 예상 소요 8분.}
\end{frame}

\begin{frame}{Must 공수 비중 — 작은 숫자로 먼저}
\begin{itemize}
\item 가용 24. 예비 25\% = 6.0 고정. Must 상한 60\% = 14.4. Should는 나머지
\item 견적을 받아 Must를 표시하니 \textbf{15.5}(64.6\%) — 초과. 항목을 빼지 않고 \textbf{정의를 줄인다}: 음성 "문장 생성" → "변환만"(1.5 → 0.5), eval 도구 1.0 → 0.5 → \textbf{14.0}(58.3\%)
\item Should 4.0 · 예비 6.0 · 합 24.0. 범위 밖: Could 12.0 · Won't 11.0
\item 상한의 용도 — "무엇을 Must에서 빼는가"가 아니라 "Must 항목의 정의를 어디까지 줄이는가"를 묻게 하는 것
\end{itemize}
\keymsg{Must ≤ 60\%는 상한이지 목표가 아니다}
\note{교재 2.2절 "작은 숫자로 먼저". DSDM MoSCoW의 60\%는 관행(권고)이지 표준 조항이 아니다. 견적표를 PM 혼자 만들거나 PM이 견적을 깎으면 6개월 뒤 "견적이 틀렸다"와 "범위가 컸다"를 구별할 수 없다. 예상 소요 6분.}
\end{frame}

\begin{frame}{Won't — "나중에"가 아니다. 재검토 조건이 측정 가능한 사건이어야 한다}
\begin{tbl}[\footnotesize]{L{12mm}L{28mm}L{66mm}L{20mm}}
B-ID & 항목 & 재검토 조건 & 시점 \\ \midrule
B-17 & 자동판정 모드 2.0 & ㉗ 합격 기준 충족 \textbf{그리고} X-04 회피 행동 < 10\% & GA 후 첫 분기 \\
B-18 & 온프렘 배포 6.0 & 유료 5개사 달성 \textbf{또는} 1차사 채널 계약 서명 — ㉕ 항목 6-1과 같은 문장 & D+360 이후 \\
B-21 & MES 연동 3.0 & 유료 고객 중 MES 보유 과반 그리고 규격 2종 이하 & 다음 사이클 \\
(규칙) & 고객 요청 항목은 이 표를 거친다 & B-20의 재발 방지 & v1.1부터 \\
(규칙) & "검토 중"이라 답하지 않는다 & 재검토 조건을 그대로 말한다 — 가격표 제외 칸이 공개판 & v1.0부터 \\
\end{tbl}
\keymsg{동보프레스는 Won't를 지키며 요구의 실체를 쪼갰다(도면 미업로드 모드 0.5) — 이 문서를 거쳤다. 태산정밀 62항목은 거치지 않았다}
\note{교재 2.3절 · 워크북 ㉖ 항목 5. ㉕와 다른 문장이면 세일즈는 유리한 쪽을 말한다. 큰 Won't(온프렘)는 지키면서 작은 Could(62항목)가 "금방 되니까" 새는 것 — 작은 것이 새는 곳이다. 예상 소요 8분.}
\end{frame}

\begin{frame}{두 번째 절단 — 손실 7.0의 다섯 계단}
\fig[0.56]{fig_2_3_capacity_loss.pdf}
\keymsg{퇴사가 원인인 것은 3.9, 3.1은 퇴사 전에 새어 나갔다 — 두 번째 계단이 조현우의 퇴사 사유와 같은 항목이다}
\note{그림 2-3 · 교재 2.4절. 조현우 잔여 2.2 / 태산 커스텀(계획 외) 2.2 / 인수인계·채용 1.0 / WIP 폐기 0.7 / B-02 조도 재작업 0.9. 분해가 없으면 대응은 "채용"뿐(Brooks — 채용은 셋째 계단을 더 먹는다), 분해가 있으면 규칙이 된다. 그 요청을 받은 사람은 CPO — ㉑ 항목 5의 거부권은 절단을 막는 권한이었지 계획 외 착수를 막는 권한이 아니었다. 정본 §1.2 "요구사항 정리자"의 첫 시험. 이미 쓴 10.7은 되돌릴 수 없다 → 남은 5.5로 Must 1.1·Should 1.5·예비 2.9. 비례 축소(24 × 17/24)는 틀린 답. 예상 소요 10분.}
\end{frame}

\begin{frame}{1교시 핵심 정리}
\begin{enumerate}
\item 시드의 결정 다섯 — 되돌리는 비용이 크다. 그래서 숫자와 서명이 많다
\item 시계 셋 — 런웨이·클로징·조건부 자금. D+149 1,505만이 바닥
\item v1.0을 사건 전에 적어야 v1.1이 사후 설명이 아니라 측정이 된다
\item Appetite는 기간 · 견적은 안 깎는다 · Must 상한은 정의 축소를 묻는 장치
\item Won't = 재검토 조건, ㉕와 같은 문장 · 손실은 분해(3.9 + 3.1) · 남은 공수로 다시 나눈다
\end{enumerate}
\keymsg{범위 시점에 잘못 적힌 것은 측정 시점에 고쳐지지 않는다}
\note{교재 제1·2장 핵심 정리 "한 줄씩" 상자의 발췌. 예상 소요 3분.}
\end{frame}

% =====================================================================
\section{2교시 (90분) — 게이트 · PR/FAQ · 사전등록 · AI 완료 정의 · 분자와 분모 · 데이터셋}
% =====================================================================
\begin{frame}{이 교시에서 배우는 것}
\begin{itemize}
\item 게이트는 날짜가 아니라 숫자와 판정자 한 사람 — 그리고 되돌림 조건 (2.5절)
\item 출시일의 글을 지금 쓴다 · 내부 FAQ의 "미확정" · 사전등록 OEC·표본·peeking (2.6\til{}2.8절)
\item 딥게이지의 6개월 — 29건·두 절단·네 실험 (2.9절)
\item AI 기능의 완료 정의는 왜 다른가 · FN·FP의 분모 · 데이터셋이 대표하는 것 (3.1\til{}3.3절)
\end{itemize}
\keymsg{Day3에서 발주자로서 요구했던 세 칸(측정·판정자·모집단)을 오늘 공급자로서 채운다}
\note{교재 2.5\til{}2.9절 + 3.1\til{}3.3절. 90분 — 2.5절 8분, 2.6절 8분, 2.7절 8분, 2.8절 12분(그림 2-4), 2.9절 6분, 3.1절 12분(그림 3-1), 3.2절 15분(그림 3-2), 3.3절 15분(그림 3-3) + 정리 6분. 예상 소요 2분.}
\end{frame}

\begin{frame}{게이트 — 진입 기준(숫자) · 판정자 한 사람 · 되돌림 조건}
\begin{tbl}[\footnotesize]{L{14mm}L{68mm}L{14mm}L{32mm}}
게이트 & 진입 기준 & 판정자 & 되돌림 \\ \midrule
알파 D+200 & 내부 + 세영 1라인 · 자동 생성 성공률 ≥ 90\% · 크래시 0건/주 · 매핑 완료 & CTO & 성공률 < 80\% 2주 → 내부로 \\
베타 D+240 & 3사 전 라인 · 기록 ≤ 25분/일 · 오전 OCR ≥ 90\% · 재촬영 < 15\% · 초안은 세영 오전 1라인만 & CPO & 기록 > 35분 또는 재촬영 > 25\% → 알파로 \\
GA D+300 & 유료 3사 · eval v1 측정 · 오류 예산 가동 · HITL · 자동판정 OFF · \textbf{미합격 시 "합격 전 베타"} & CEO & 예산 2개월 연속 소진 또는 2사 해지 → 베타로 \\
\end{tbl}
\keymsg{㉖ 항목 7과 ㉙ 항목 5에 같은 숫자로 — 한 곳에서 복사한다. 두 곳에 쓰면 갈린다}
\note{교재 2.5절. GA 기준의 "합격 전 베타"가 3교시 FN 8.1\% 판정을 가능하게 한 미리 열어 둔 길. 되돌림이 없는 GA는 일방통행 — "이미 출시했으니". PM ⑯ SAC·Day4 2.2절 go/no-go가 발주자 쪽 거울. 예상 소요 6분.}
\end{frame}

\begin{frame}{Working Backwards — 출시일의 글을 지금 쓰고, 막히는 문장을 지우지 않는다}
\begin{itemize}
\item 제목 "안산 세영정공, 검사기록 작성 하루 47분 → 19분" — 고객과 숫자. 해법 칸은 기능 목록이 아니라 검사원의 하루
\item D+158에 참이 아닌 두 문장 — "부적합 통보 1영업일 이내"(X-01은 기록 시간만 쟀다), "3개사 도입"(유료 0)
\item 지우지 않고 내부 FAQ 15번에 적었다 — D+300 전에 참이 되거나 문장이 빠진다
\item 요점은 글이 아니라 \textbf{막히는 곳} — ㉖에서 잘라 낸 것이 보도자료에 필요하면 잘못 잘랐거나 과장이다
\end{itemize}
\keymsg{보도자료를 마케팅에 넘기지 않는다 — 기획 문서다}
\note{교재 2.6절 · 워크북 ㉙ 항목 1·3. Bryar \& Carr(2021) 4장. PM 트랙에는 대응물이 없다(요구사항이 주어지므로) — 가장 가까운 것이 Day1 비즈니스 케이스의 편익 문장. 예상 소요 6분.}
\end{frame}

\begin{frame}{내부 FAQ 15문 — 미확정 넷이 정상이다}
\begin{tbl}[\footnotesize]{L{6mm}L{52mm}L{68mm}}
\# & 질문 & 답 \\ \midrule
4 & 단가 45만은 원가를 감당하는가? & 표준 13.4\%, 헤비 51.5\%. \textbf{헤비가 몇 곳인지는 미확정 — 계측 없음} \\
7 & 검사원이 초안을 받아들이는가? & \textbf{미확정 — X-04}(D+290). X-01에서 5명 중 2명 별도 종이 기록 \\
8 & 공장이 라인 과금을 받아들이는가? & \textbf{미확정 — X-05} \\
9 & 조도 경고가 재촬영을 줄이는가? & \textbf{미확정 — X-07}(D+280) \\
10 & 유료 5개사(D+360)는 가능한가? & \textbf{미확정 — 관찰만.} 실험이 아니라 세일즈 \\
11 & FN 8.1\%(미합격)로 GA해도 되는가? & 기준 유지, 배포 조건을 낮춘다 — "합격 전 베타" \\
15 & 이 보도자료에서 아직 참이 아닌 문장은? & "통보 1영업일" · "3개사 도입" \\
\end{tbl}
\keymsg{15문 전부에 답이 있으면 실패 — 어려운 질문을 뺐거나 아는 척한 것이다}
\note{교재 2.7절. 미확정 넷은 ㉖ 항목 8의 리스크 다섯에서 왔다 — 실험이 있으면 X-nn, 없으면 "관찰만". 투자자용으로 다듬는 순간 미확정이 사라진다. 4번의 "계측 없음"이 5교시 청구서로 돌아온다. 예상 소요 6분.}
\end{frame}

\begin{frame}{사전등록 — OEC 하나 · 표본과 근거 · 분석 시점 고정 · peeking 금지 · SRM}
\fig[0.55]{fig_2_4_sample_size.pdf}
\keymsg{효과가 작을수록 표본은 제곱으로 는다 — 사전등록은 "무엇을 검출할 수 없는가"를 먼저 적는 일이다}
\note{그림 2-4 · 교재 2.8절 "작은 숫자로 먼저". X-07: 18\% → 8\%, 두 비율 검정, α .05 양측·검정력 .8 → 177/arm. 12\%면 약 500, 14\%면 약 1,100 — 2주에 얻는 250으로는 12\% 이하 효과를 검출할 수 없다. X-04의 OEC는 "수용률 ≥ 70\% 그리고 회피 < 10\%" — Day5 X-01의 교훈을 문장으로. X-05·X-06은 모집단 3 → "전수", 성공 3/3. 중단 기준(회피 20\% → 초안 OFF)은 가드레일. 예상 소요 10분.}
\end{frame}

\begin{frame}{딥게이지의 6개월 — 29건 · 두 절단 · 네 실험}
\begin{itemize}
\item \textbf{D+130} 29건 41.0 → Must 15.5 → 14.0(정의 축소) · Should 4 · 예비 6 = 24. Won't 셋에 재검토 조건. 게이트 셋. 리스크 5 → 실험 4 + 관찰 1
\item \textbf{D+158} 보도자료 8칸 · 고객 FAQ 12 · 내부 FAQ 15(미확정 4) · 사전등록 4 서명
\item \textbf{D+200} 동보프레스 도면 미업로드 모드(B-19, 0.5, 예비) — 이 문서를 거쳤다 / \textbf{D+180\til{}228} 태산정밀 62항목(B-20) 2.2 — 거치지 않았다 / \textbf{D+230} 퇴사
\item \textbf{D+235} 손실 7.0 분해 → 남은 5.5: B-08 내보내기 0.4 · B-10 D1\til{}D3 0.5 · B-25 CTO 직접 · B-22 → Won't · B-14 세영 1라인 1.0 · 예비 2.9
\item \textbf{D+280·290} X-05 성공(3/3) · X-06 성공(9·6·10일) · X-07 \textbf{불확정}(9.4\%) · X-04 \textbf{불확정}(수용률 71\% / 회피 14\%) → 둘 다 미리 적힌 처리대로
\end{itemize}
\keymsg{넷 중 둘이 불확정 — 그리고 둘 다 사전등록에 불확정의 처리가 미리 적혀 있었다}
\note{교재 2.9절. 학생이 4교시에 쓸 ㉖의 사건 흐름을 한 장으로. 이 교재는 D+300에 멈춘다 — 그 뒤는 Day7 워크북 §0.3. 예상 소요 6분.}
\end{frame}

\begin{frame}{AI 기능의 완료 정의 — 같은 표의 반대편}
\fig[0.58]{fig_3_1_rules_of_credit.pdf}
\keymsg{"측정·판정자·모집단 없는 완료는 완료가 아니다" — Day3에 벤더에게 한 말을 오늘 우리가 쓴다}
\note{그림 3-1 · 교재 3.1절. 일반 기능은 "동작한다", AI 판정 기능은 "어느 데이터에서 얼마나 맞는가". 세 칸 + 배포 뒤의 둘(오류 예산·HITL) = ㉗ 한 묶음. ㉖ 항목 6의 AI DoD는 ㉗ 항목 3 인용만 — 두 곳에 쓰면 갈린다. 기준(D+151)이 측정(D+270)보다 넉 달 앞선다는 순서. 예상 소요 10분.}
\end{frame}

\begin{frame}{분자와 분모 — FN 24/296, FP 173/1,184. 분모가 다르다}
\fig[0.58]{fig_3_2_confusion.pdf}
\keymsg{판정 기능에 "정확도"는 쓰지 않는다 — 전부 양품이라 해도 80\%가 나온다. 위험한 칸은 FN}
\note{그림 3-2 · 교재 3.2절 "작은 숫자로 먼저". 판정 단위(항목) 먼저 — 60건 = 1,480항목. FN = 실제 불량 중 초안 양품(안전), FP = 실제 양품 중 초안 불량(수용 — FP가 늘면 검사원이 초안을 무시한다). OCR 960/1,050 = 91.4, 외관 F1 = 2·0.768·0.792/(0.768+0.792) = 0.78. 24 + 173 = 197은 "오류 197건"이 아니다. FN을 낮추면 FP가 오르는 긴장은 Day7. 예상 소요 12분.}
\end{frame}

\begin{frame}{데이터셋이 대표하는 것 — 그리고 대표하지 않는 것}
\fig[0.58]{fig_3_3_dataset.pdf}
\keymsg{여집합의 첫 줄(오후 조도)은 지금 일한다. 나머지 줄은 지금 아무 일도 하지 않는다 — 비우지 않는 습관만 가져가라}
\note{그림 3-3 · 교재 3.3절. "3세트 60건"은 구성이지 모집단이 아니다. 모집단 문장 한 줄 + 여집합 다섯 줄 — 고객이 실제로 가진 것부터. 라벨러(전 검사원 2인 + 강민수 검수, 불일치 6.8\%)도 모집단의 일부 — 정답의 불확실성 위에 FN 8.1이 있다. Day3의 91.3\% 분모(순환실사 SKU 69\%)와 같은 자리 — 그때는 감사자, 지금은 쓴 사람. 나머지 줄이 무엇을 하는지는 Day8 — 지금은 말하지 않는다(F-11). 예상 소요 12분.}
\end{frame}

\begin{frame}{2교시 핵심 정리}
\begin{enumerate}
\item 게이트 = 숫자 + 판정자 한 사람 + 되돌림. ㉖·㉙에 같은 숫자로
\item 출시일의 글을 지금 쓰고 막히는 문장을 지우지 않는다 · 내부 FAQ 15문 전부에 답이 있으면 실패
\item 사전등록 — OEC 하나, 표본과 근거(또는 전수), 분석 시점 고정, peeking 금지, SRM
\item AI 완료 정의 = 측정·판정자·모집단 세 칸 + 오류 예산 + HITL — 발주자로서 요구했던 그 칸
\item FN과 FP의 분모는 다르다 · 정확도는 쓰지 않는다 · 모집단 문장과 여집합을 적는다
\end{enumerate}
\keymsg{결과를 보기 전에 기준을 — 실험에서도, eval에서도}
\note{교재 제2·3장 핵심 정리 발췌. 예상 소요 3분.}
\end{frame}

% =====================================================================
\section{3교시 (90분) — 합격 기준·개정 · 오류 예산 · HITL · 밸류메트릭 · UE 두 열}
% =====================================================================
\begin{frame}{이 교시에서 배우는 것}
\begin{itemize}
\item 합격 기준은 데이터셋 버전과 함께 · 개정 절차를 같은 날 · FN 8.1\%의 불합격 (3.4절)
\item 오류 예산 — 측정 원천이 먼저, 놓치는 것을 적고, 소진 시 조치는 자동 (3.5절)
\item HITL — 자동판정을 켜지 않는 설계 · 문구 둘 · 현장 수용성 (3.6절) · 분야별 대조 (3.7절, 심화)
\item 밸류메트릭 4안 × 5기준 · 45만의 근거 · 유닛이코노믹스 두 열 (4.1\til{}4.2절)
\end{itemize}
\keymsg{기준과 배포는 다른 결정이다 — 기준을 지키고 배포 조건을 낮춘다}
\note{교재 3.4\til{}3.7절 + 4.1\til{}4.2절. 90분 — 3.4절 15분, 3.5절 15분(그림 3-4), 3.6절 15분(그림 3-5), 3.7절 5분, 4.1절 15분(그림 4-1), 4.2절 15분(그림 4-2) + 정리 5분 + 예비 5분. 예상 소요 2분.}
\end{frame}

\begin{frame}{합격 기준과 개정 절차 — criteria drift를 결과 전에 막는다}
\begin{itemize}
\item 기준은 "FN ≤ 8.0\% 그리고 OCR ≥ 90\% \textbf{on DS-v1}" — 데이터셋 버전이 없는 기준은 기준이 아니다
\item 개정 절차를 기준과 같은 날: \textbf{낮추는} 개정(FN 상향·OCR 하향·DS 축소) = 대표 승인 + 이사회 보고 + 외부 근거(측정 오류의 입증 또는 모집단 변경) / \textbf{높이는} 개정 = CTO 전결
\item 사고·일정·영업 압력은 근거가 아니다 — 그것을 적는 유일한 시점이 결과 전이다
\item Shankar 외(2024): 평가 기준은 결과를 보면서 바뀐다 — 8.1이 나오면 8.1로, 8.3이 나오면 8.3으로
\end{itemize}
\keymsg{개정 절차가 없는 합격 기준은 결과 뒤에 결과에 맞춰진다}
\note{교재 3.4절. Nosek 외(2018) 사전등록 논리와 같은 자리. PM Day3 "Rules of Credit은 측정이 아니라 합의의 산물"(2.2절) — 합의를 결과 뒤에 바꾸면 측정이 아니다. 예상 소요 7분.}
\end{frame}

\begin{frame}{FN 8.1\% — 0.1\%p, 24건 대 23건. 불합격이다}
\begin{itemize}
\item D+270 측정: OCR 91.4\% ✓ · \textbf{FN 8.1\%}(기준 8.0) · FP 14.6\% · F1 0.78 · 1.8초 · 32원
\item "사실상 합격"이라고 적고 싶어진다 — 그러나 한 건 차이는 측정 오류의 입증도 모집단 변경도 아니다
\item 예시본의 판정: \textbf{불합격, 기준 유지}. 그리고 배포 조건을 낮춘다 — "합격 전 베타", 세영 오전 1라인, 승인 필수, 자동판정 없음
\item 그 길은 ㉖ 항목 7 GA 기준에 미리 열려 있었다 — "초안 미합격 시 '합격 전 베타' 표시"
\end{itemize}
\keymsg{순서 — 기준(D+151) → 측정(D+270) → 판정(유지) → 배포 조건(낮춤). 거꾸로 하면 모든 모델이 합격한다}
\note{교재 3.4절 딥게이지 사례. 이 판정이 옳았는지는 이 교재가 답하지 않는다 — Day7. 학생이 4교시 ㉗ 항목 3에서 8.1\%를 받았을 때 무엇을 하는지 쓰게 한다 — 8.1로 고친 조는 즉시 탈락. 예상 소요 8분.}
\end{frame}

\begin{frame}{오류 예산 — 측정 원천은 HITL 수정 로그, 놓치는 것은 클레임}
\fig[0.55]{fig_3_4_error_budget.pdf}
\keymsg{예산 8.0\%는 합격 기준과 같은 숫자이지만 분모가 다르다 — 검사원도 놓친 FN은 이 그림에 없다}
\note{그림 3-4 · 교재 3.5절. SRE error budget의 번역 — 그러나 배포 뒤 FN은 라벨이 없다. 측정 원천 = 검사원이 초안을 뒤집은 기록 → "검사원이 잡은 FN"이지 실제 FN이 아니다. 소진 시 조치는 자동(라인 초안 OFF → 재학습 → 재측정 → ON, 릴리스 동결). 신뢰도 0.85 미만 초안 미표시 — 예산을 지키는 가장 싼 방법. 그림 값은 가상 시나리오. 예상 소요 10분.}
\end{frame}

\begin{frame}{HITL — 촬영에서 로그까지, 자동판정은 만들지 않았다}
\fig[0.58]{fig_3_5_hitl_flow.pdf}
\keymsg{초안 카드는 세 조건의 AND — 베타 라인 · 신뢰도 ≥ 0.85 · 오류 예산 미소진. 아니면 "판정 보류"}
\note{그림 3-5 · 교재 3.6절. 네 칸(ON/OFF 조건·승인 흐름·에스컬레이션·문구) + 현장 수용성. 수정 사유 1탭과 "수정 건수 = 모델 품질 지표" 표기가 있어야 수정 로그가 살고 그래야 오류 예산이 산다 — 두 항목이 서로에게 달려 있다. ISO/IEC 42001 인간 감독. 예상 소요 8분.}
\end{frame}

\begin{frame}{문구 둘, 그리고 박선재의 눈}
\begin{itemize}
\item \textbf{화면}(검사원에게): "판정 초안은 참고용입니다. 최종 판정은 검사원이 합니다." — 판정권을 \textbf{돌려주는} 문장
\item \textbf{계약}(고객사에게, ㉘ 항목 6): "당사 책임은 연 계약액의 10\%(162만)를 상한으로" — 상한을 \textbf{알리는} 문장
\item 같은 화면에 두면 검사원은 "나머지 90\%는 내 책임"으로 읽는다 — 촉탁직에게 그것은 재촬영과 별도 종이 기록으로 돌아온다(X-01)
\item 현장 수용성 칸: A-05b \textbf{미해결로 표기}. 대응은 교육이 아니라 설계 — 수정이 벌점이 아니어야 수정 로그가 산다
\end{itemize}
\keymsg{문구가 두려움을 없애지는 않는다 — 그래도 두 문장을 같은 화면에 두지 않는다}
\note{교재 3.6절 · 워크북 ㉗ 항목 6. 정본 §8.3 "자동판정을 끈다 — 고용 구조의 반응"은 D+300 이후. 지금 한 것은 자동판정을 만들지 않은 것과 초안을 세영 오전 1라인에만 켠 것. X-04 결과(71\% / 14\% 불확정)가 그 결정을 반쯤 지지한다. 예상 소요 7분.}
\end{frame}

\begin{frame}{분야별 대조 — 사람이 판정을 뒤집을 수 있는가 (심화)}
\begin{tbl}[\footnotesize]{L{22mm}L{28mm}L{38mm}L{38mm}}
분야 & FN의 뜻 & 사람의 자리 & 딥게이지와의 거리 \\ \midrule
제조 검사(딥게이지) & 불량이 라인을 나간다 & 검사원 승인 필수, 자동판정 없음 & — \\
의료 영상 보조 & 놓친 병변 & 판독의가 최종, "보조" 등급 & 가장 가깝다 \\
금융 사기 탐지 & 사기 통과 vs 정상 차단 & 임계값 위 자동, 중간만 사람 & 자동이 기본 — 반대 \\
자율주행 & 미감지 = 사고 & 운전자 감독, 개입 요구 & 반응 시간이 설계 \\
콘텐츠 검수 & 위반 통과 vs 정상 삭제 & 자동 삭제 + 이의 제기 & 대량·저비용 — 반대 \\
\end{tbl}
\keymsg{딥게이지는 의료 쪽 선택을 했다 — 대가는 "AI 검사"라는 말을 마케팅에 못 쓰는 것(㉕ 항목 6-3)}
\note{교재 3.7절(심화). 공통점은 FN·FP 비용의 비대칭, 차이는 사람이 뒤집는 것이 정상인가 예외인가. 예상 소요 4분.}
\end{frame}

\begin{frame}{밸류메트릭 4안 × 5기준 — 라인은 원가 정렬 M을 알고 골랐다}
\fig[0.58]{fig_4_1_value_metric.pdf}
\keymsg{종합은 합산이 아니라 "어느 기준을 포기했는가"의 문장 — 알고 포기했으면 계측을 먼저 넣었어야 했다}
\note{그림 4-1 · 교재 4.1절. 다섯 기준의 무게는 단계마다 다르다 — 첫 유료 3사 6개월에는 팔기 쉬운가·예측 가능한가가 무겁고 원가 정렬은 가볍다(고객 셋, 원가 210만). 좌석 L의 근거는 커널 행동 2 — 좌석 과금은 "계정 아끼기"를 만든다. 딥게이지는 알고도 계측을 안 넣었다(B-24에 사용량 계측 없음) — 그것이 5교시 780만. 선택의 대가(점포 등)는 오늘 말하지 않는다. 예상 소요 8분.}
\end{frame}

\begin{frame}{45만의 근거 — 팀장에게 하는 말}
\begin{itemize}
\item 라인당 전기·중복 입력의 돈 환산 \textbf{39.3만/월}(Day5) + 8D 11.5시간 + 통보 지연 페널티 + 클레임 회피
\item 클레임 1건 2,700만 = 라인 3.0 × 45만 × \textbf{20개월} — 한 건을 막으면 20개월치
\item 좌석 무제한 — 검사원이 늘어도 비용이 안 늘어야 검사원이 쓴다(A-01·A-05). 라인은 공장이 이미 세는 단위(X-05: 3/3)
\item 근거는 팀장에게 하는 말이지 검사원에게 하는 말이 아니다(A-03 — 팀장이 돈을 낸다)
\end{itemize}
\keymsg{선택의 대가 칸을 비우지 않는다 — 사진 수가 4.7배인 고객에서 변동원가율이 13.4 → 51.5\%}
\note{교재 4.1절 딥게이지 사례 · 워크북 ㉘ 항목 2. 밸류메트릭을 바꾸는 비용(모든 계약·청구·보고)을 모르고 "나중에 바꾸면 된다"고 하는 것이 흔한 오류. 예상 소요 5분.}
\end{frame}

\begin{frame}{유닛이코노믹스 v1 — 같은 매출 135만, 공헌이익률 86.6 vs 48.5\%}
\fig[0.58]{fig_4_2_ue_two.pdf}
\keymsg{한 고객의 산수를 두 번 — 헤비 열을 "예외"로 지우면 표준 열은 평균이 아니라 희망이 된다}
\note{그림 4-2 · 교재 4.2절 "작은 숫자로 먼저". 4,100 × 32원 = 13.1만 + 3.8 + 1.2 = 18.1 / 19,400 × 32 = 62.1 + 3.8 + 3.6 = 69.5. 공헌이익 0의 사진 수 약 39,900장 — 헤비의 두 배. 고정비(GPU·인건비)는 여기 없다 — 회사 GM에. NDR: 라인 과금은 공장 라인 수(3.2)에 묶인다 → +0.2, 실사용 2.4/3.0은 축소 위험이 먼저 → 상한 약 105\%, GRR 먼저. 예상 소요 10분.}
\end{frame}

\begin{frame}{3교시 핵심 정리}
\begin{enumerate}
\item 기준은 데이터셋 버전과 함께, 개정 절차는 같은 날 — 낮추는 개정에 외부 근거와 상위 승인
\item FN 8.1\%는 불합격 — 기준을 지키고 배포 조건을 낮춘다("합격 전 베타")
\item 오류 예산의 첫 칸은 측정 원천 — 검사원도 놓친 FN은 클레임에 있다. 소진 시 조치는 자동
\item HITL — 자동판정 없음, 세 조건의 AND, 문구 둘, 현장 수용성은 촉탁직의 눈으로
\item 밸류메트릭은 4안 × 5기준, 대가 칸을 비우지 않는다 · UE는 두 열, 고정비는 섞지 않는다
\end{enumerate}
\keymsg{오류 예산과 HITL은 서로에게 달려 있다 — 수정이 벌점이면 예산은 늘 남는다}
\note{교재 제3·4장 핵심 정리 발췌. 예상 소요 3분.}
\end{frame}

% =====================================================================
\section{4교시 (90분) — 실습 ①: ㉖ MVP 범위 결정서 + ㉗ Eval 명세서·오류 예산·HITL}
% =====================================================================
\begin{frame}{이 교시에서 하는 것}
\begin{itemize}
\item \textbf{㉖ MVP 범위 결정서} — 워크북 §1.4에 항목 3(산식 두 번) · 4(절단 — 41 → 24는 MoSCoW 열로, 24 → 17은 손실 분해와 재절단으로) · 5(Won't) — \textbf{40분}
\item \textbf{㉗ Eval 명세·오류 예산·HITL} — 워크북 §2.4에 항목 2(데이터셋 — 모집단·여집합) · 3(합격 기준·개정 절차) · 5(오류 예산) · 6(HITL — 문구 둘) — \textbf{45분}
\item 순서 — 읽는다(X.1\til{}X.3) → 찾는다(◆ 카드: 백로그 29건 견적표, D+235 상태, 데이터셋 분자·분모) → 쓴다(빈 템플릿) → \textbf{대조한다}(X.5·X.6은 다 쓴 뒤)
\item 오늘은 카드에 찾을 것이 많다 — 카드에 없는 숫자는 만들지 않는다
\end{itemize}
\keymsg{정답을 보지 말고 쓰고, 다 쓴 뒤 대조하라 — 해설은 정답이 아니라 즉시 탈락 조건과 방어 가능한 범위다}
\note{워크북 §0.1 따라가는 순서. 40 + 45 + 안내·정리 5 = 90. 조별 4\til{}5인. 인쇄용 빈 템플릿(00\_Day6\_템플릿팩\_빈양식.pdf 20쪽 — 견적표·루브릭·사전등록서는 가로)을 미리 배부. 예상 소요 3분.}
\end{frame}

\begin{frame}{㉖ MVP 범위 결정서 — 과제 지시}
\begin{itemize}
\item \textbf{과제}: 항목 3(Must 산식 — 24 기준과 17 기준) · 4(절단 결정) · 5(Won't 명시 목록)를 딥게이지로 채우라. 항목 2의 MoSCoW 열이 항목 3·4의 입력이므로 먼저 표시한다
\item 시점 v1.0 \textbf{D+130} / v1.1 \textbf{D+235}. 카드가 주는 것 — 백로그 29건 41.0 mm(B-ID·항목·상위 노드·mm), 가용 24 = 4명 × 6개월, D+235 상태(완료 10.0 · WIP 0.7 · 예비 사용 0.8 · \textbf{태산 커스텀 2.2} · 재작업 0.9 · WIP 폐기 0.7 · 채용 1.0 · 잔여 2.2), TIPS 조건
\item 앞 산출물 — ㉓ 항목 4·7(우선순위·기각), ㉕ 항목 6(하지 않는 것), ㉔ 미제거 가정 4
\item 여러분이 결정하는 것 — MoSCoW 배분 · 상한 규칙 · 24 → 17에서 무엇을 축소·이동 · 재검토 조건 · 통보 대상
\end{itemize}
\keymsg{상위 노드가 트리의 잎이 아닌 6건이 절단의 첫 질문 — 어디서 왔는가}
\note{워크북 §1.1\til{}1.3 + ◆ 카드. 교재 2.1\til{}2.4절은 1교시에 배웠다. 견적표는 주어진다 — 학생이 하는 것은 MoSCoW 열·산식·절단·Won't. 손실 분해 5줄의 재료는 카드에 있으므로 "퇴사 7.0" 한 줄로 적은 조는 즉시 탈락. 예상 소요 5분.}
\end{frame}

\begin{frame}{작성 요령 ㉖ — 산식 두 번 · 손실 분해 · Won't (항목 3·4·5)}
\begin{itemize}
\item \textbf{항목 3}: 상한 규칙을 먼저(예: Must ≤ 60\%, 예비 25\%, Should ≤ 15\%), 그 다음 계산. 17 기준에서는 \textbf{쓴 것과 남은 것을 분리} — 비례 축소는 틀린 답
\item \textbf{항목 4}: 41 → 24는 Could·Won't로 보낸 것과 정의를 줄인 것(15.5 → 14.0)을 적는다. 24 → 17은 손실 7.0을 다섯 줄로 분해한 뒤 남은 5.5로 무엇을 끝내는지 — 축소는 바뀐 정의를 적는다("자동 제출 → 내보내기")
\item \textbf{항목 5}: 재검토 조건은 측정 가능한 사건 — "고객이 원하면"은 조건이 아니다. ㉕ 항목 6과 같은 문장. B-20의 재발 방지 규칙 한 줄
\item 통보 대상 열을 비우지 않는다 — 그 기능이 온다고 믿는 고객·투자자·세일즈
\end{itemize}
\keymsg{절단의 기술 — 항목을 빼지 않고 정의를 줄인다}
\note{워크북 §1.3 항목 3·4·5 가이드 + §1.6 판정 규칙. 방어 가능한 다른 답 — B-14(판정 초안 베타)를 Could로 내린 것(근거: ㉗ 미합격), 예비 20\%, Must 상한 65\%. 예상 소요 5분.}
\end{frame}

\begin{frame}{㉖ 흔한 오류 점검}
\begin{itemize}
\item Appetite 칸에 기능 목록이나 "24 mm"를 적었다 — 기간이 아니다 → 탈락
\item 17 mm 기준을 24 × 17/24로 비례 축소했다 → 탈락. 이미 쓴 10.7은 비례하지 않는다
\item Won't 칸에 재검토 조건 없이 "Later"·"검토 중" → 탈락
\item 24 → 17의 손실을 "퇴사 7.0" 한 줄로 → 탈락. 다섯 줄 중 세 줄만 퇴사가 원인이다
\item 정답 처리: MoSCoW 배분이 예시본과 달라도 상한 규칙을 세우고 지켰고 절단마다 사유·통보 대상이 있는 것
\end{itemize}
\keymsg{손실 분해가 없으면 대응은 "채용"뿐이고, 분해가 있으면 대응은 규칙이 된다}
\note{워크북 §1.6 판정 규칙. 조별 순회 시 항목 4의 손실 분해 표를 먼저 본다. 예상 소요 4분.}
\end{frame}

\begin{frame}{㉗ Eval 명세서·오류 예산·HITL — 과제 지시}
\begin{itemize}
\item \textbf{과제}: 항목 2(데이터셋 — 모집단 문장과 "대표하지 않는 것") · 3(합격 기준·개정 절차 — 그리고 FN 8.1\%를 받았을 때 무엇을 하는가) · 5(오류 예산) · 6(HITL — 화면 문구와 계약 문구)
\item 시점 \textbf{D+151}(측정 기록 D+270). 카드가 주는 것 — 정본 V-01\til{}07 v1 열, 데이터셋 내역(3세트 × 20건, 오전 9\til{}11시, 1,480항목 = 불량 296 / 양품 1,184, 치수 1,050, 외관 430), 측정치 분자(960 · 24 · 173 · 76/20/23), HITL 수정 로그의 존재, 책임 한계 162만
\item 앞 산출물 — ㉔ X-01(수용 실패)·X-02(조도), ㉖ 항목 6 DoD, PM ⑫ Rules of Credit 세 칸, PM ⑰ error budget 형식
\item 여러분이 결정하는 것 — 모집단 문장·여집합 · 개정 절차 · 예산 비율·임계값·조치 · 문구 둘 · 현장 수용성 한 줄
\end{itemize}
\keymsg{항목 2의 한 문장(어느 산업·공정·조도·몇 건)이 Day8 마지막 30분까지 따라간다 — 정확히 적기만 하라}
\note{워크북 §2.1\til{}2.3 + ◆ 카드. 교재 3.1\til{}3.6절은 2·3교시. "자동차 3세트"를 학생이 쓰게 하되 그 함의는 말하지 않는다(F-11, 제작 노트 §4 Day6). 예상 소요 5분.}
\end{frame}

\begin{frame}{작성 요령 ㉗ — 모집단·개정 절차·오류 예산·문구 둘 (항목 2·3·5·6)}
\begin{itemize}
\item \textbf{항목 2}: 모집단 = 산업·공정·조명·촬영 조건 한 문장. 여집합 = 그 문장의 여집합 중 \textbf{고객이 실제로 가진 것}부터(오후 조도). 라벨러가 개발자면 안 된다
\item \textbf{항목 3}: "FN ≤ 8.0\% on DS-v1". 낮추는 개정과 높이는 개정을 분리. 8.1\%를 받으면 — 기준 유지, 배포 조건 낮춤. 8.1로 고친 조는 탈락
\item \textbf{항목 5}: 측정 원천을 먼저(HITL 수정 로그). 놓치는 것을 근거 열에("검사원도 놓친 FN은 클레임으로"). 소진 시 조치는 자동
\item \textbf{항목 6}: 자동판정 ON 조건이 "신뢰도 ≥ n"뿐이면 탈락. 화면 문구와 계약 문구를 분리. 현장 수용성은 박선재의 눈으로 — "교육으로 해결"은 답이 아니다
\end{itemize}
\keymsg{탈락 조건 — 여집합 공란 · FN·FP 같은 분모 · 8.1로 고침 · 문구 하나}
\note{워크북 §2.3 가이드 + §2.6 판정 규칙. 정답 처리 — 예산 비율·임계값이 달라도 측정 원천과 자동 조치가 있으면; "합격 전 베타" 대신 "GA 연기"를 택했으면(근거가 있으면). 예상 소요 5분.}
\end{frame}

\begin{frame}{㉗ 흔한 오류 점검 — 그리고 두 문서가 만나는 곳}
\begin{itemize}
\item 항목 1에서 FN·FP의 분모를 같은 것으로 적었다 → 탈락(불량 296 / 양품 1,184)
\item 오류 예산 FN 8.0\%와 합격 기준 FN 8.0\%가 같은 것을 잰다고 적었다 → 분모가 다르다(실제 불량 vs 검사원이 뒤집은 것)
\item 계약 문구가 검사원 화면에 뜬다 → "나머지 90\%는 내 책임". 촉탁직의 재촬영·별도 기록
\item 만나는 곳 — ㉗ 항목 6 계약 문구 = ㉘ 항목 6(162만, 한 곳에서 복사) / ㉗ 오류 예산이 놓치는 FN = ㉘ 책임 한계가 맡는 클레임 / ㉖ 항목 6 AI DoD = ㉗ 항목 3 인용
\end{itemize}
\keymsg{같은 구멍을 두 문서가 반대편에서 보고 있다 — 예산이 놓치는 FN과 책임 한계 162만}
\note{워크북 §2.6 해설. 예상 소요 4분.}
\end{frame}

\begin{frame}{예시본 참조 — 워크북 §1.5\til{}1.6 · §2.5\til{}2.6}
\begin{itemize}
\item \textbf{㉖} v1.1 — Must 14.0(58.3\%) → 17 기준 11.8(69.4\%, 남은 5.5로 Must 1.1·Should 1.5·예비 2.9) · 손실 5줄 · 절단 8행 · Won't 4 + 규칙 2 · 게이트 3 · 리스크 5 → 실험 4 + 관찰
\item \textbf{㉗} v1.0 — 판정 단위 항목 · DS-v1 모집단 문장 + 여집합 5줄 · 기준 v1 + 낮추는/높이는 개정 분리 · 측정 기록 D+270(불합격, 기준 유지) · 루브릭 4 × 4 · 오류 예산(HITL 로그, 0.85, 자동 OFF) · HITL 문구 둘 + A-05b 미해결
\item 대조 규칙 — 칸별로. 다른 칸은 §X.6에서 즉시 탈락인지 방어 가능한지 확인. 예시본과 다른 것이 오답이 아니다
\end{itemize}
\keymsg{인쇄용 완성 예시 — 00\_Day6\_템플릿팩\_완성예시.pdf(15쪽)}
\note{예시본은 저녁에 읽는 것이 정상. 수업 중에는 판정 규칙 상자만 같이 본다. 예상 소요 3분.}
\end{frame}

\begin{frame}{4교시 정리 — 두 문서에서 반복된 규칙}
\begin{itemize}
\item \textbf{정의를 줄인다} — Must 15.5 → 14.0, 포털 자동 제출 → 내보내기 / 기준을 지키고 배포 조건을 낮춘다
\item \textbf{분해한다} — 손실 7.0을 다섯 줄로 / FN·FP를 다른 분모로 / 모집단과 여집합으로
\item \textbf{한 곳에서 복사한다} — 게이트 숫자(㉖ ↔ ㉙), AI DoD(㉖ → ㉗), 계약 문구(㉗ → ㉘), Won't 문장(㉕ → ㉖)
\item \textbf{놓치는 것을 적는다} — 여집합 · 예산이 못 잡는 FN · 통보 대상
\end{itemize}
\keymsg{두 곳에 쓰면 갈린다 — 오늘 다섯 문서가 서로를 인용하는 이유}
\note{4교시 마무리. 예상 소요 3분.}
\end{frame}

% =====================================================================
\section{5교시 (75분) — 청구서 · ARR · 계약 다섯 칸 · 텀시트 · 캡테이블 · 참가적 · 옵션풀}
% =====================================================================
\begin{frame}{이 교시에서 배우는 것}
\begin{itemize}
\item 청구서 780만이 말하는 것과 말하지 않는 것 — 가설 넷, 계측 (4.3절)
\item 구축비와 ARR — 4,860 vs 6,060, 같은 원가의 다른 분모 (4.4절) · 계약의 다섯 칸과 가격표의 제외 칸 (4.5절) · 분야별 대조 (4.6절, 심화)
\item 텀시트를 읽는 법 — "우리에게 무슨 뜻인가" · 캡테이블은 주식수로 · 참가적 1x · 옵션풀의 위치 · TIPS 압력 · 협상과 BATNA (5.1\til{}5.6절)
\end{itemize}
\keymsg{정답이 없는 절이 하나 있다 — 청구서. "정답을 아직 알 수 없다"가 정답이다}
\note{교재 4.3\til{}4.6절 + 제5장. 75분 — 4.3절 12분(그림 4-3), 4.4절 8분(그림 4-4), 4.5절 6분, 4.6절 3분, 5.1절 8분, 5.2절 8분(그림 5-1), 5.3절 10분(그림 5-2), 5.4절 8분(그림 5-3), 5.5\til{}5.6절 8분 + 정리 4분. 예상 소요 2분.}
\end{frame}

\begin{frame}{청구서 780만 — 서비스별로는 보이고, 원인은 보이지 않는다}
\fig[0.56]{fig_4_3_cloud_bill.pdf}
\keymsg{D+284의 딥게이지는 원인을 모른다 — 정본을 읽은 여러분은 안다. 그래서 이 항목이 어렵다}
\note{그림 4-3 · 교재 4.3절. 예상 210 → 실제 780(VLM 612·STT 41·저장 58·GPU 69). 초과 570 중 동보 444(78\%) = 원가표 69.5의 6.39배(F-09). 테넌트별 합계는 있고 사진 단위 재촬영·재시도·배치 구분은 없다. 가설 넷의 대응이 전부 다르다(UX/버그/배치 정책/조항) — 어느 것인지 모르면 하나를 고르는 것은 도박. 예상 소요 8분.}
\end{frame}

\begin{frame}{정답이 없는 항목 — 가설 넷, 판별 계측, 첫 조치는 계측}
\begin{itemize}
\item 워크북 ㉘ 항목 4에 "재촬영이니 교육한다" 또는 "가격을 올린다"고 쓰면 — 정본을 읽은 사람의 답이지 D+284에 쓸 수 있는 답이 아니다
\item 방어 가능한 범위: 가설 셋 이상 + 가설마다 판별 계측 + 첫 조치 = 계측(D+300 이후 첫 백로그 항목 0.5 mm)
\item 판별 전에도 할 수 있는 것은 함께 — 공정 사용 조항(라인당 월 6,000건, 초과 건당 40원)·경고 임계값(4,500건). 단 신규·갱신부터, \textbf{소급하지 않는다}
\item 상한은 두지 않는다 — 검사원이 촬영을 아끼면 FN이 는다(㉗). 두 문서가 여기서 만난다
\end{itemize}
\keymsg{4.1절로 돌아가면 — 원가 정렬 M을 알고 골랐으면 계측을 먼저 넣었어야 했다. 알고도 안 넣은 것이 780만이다}
\note{교재 4.3절 · 워크북 ㉘ 항목 4·§3.6. 청구서를 CTO의 문제로 넘기는 것이 흔한 오류 — 가격의 문제이기도 하다. 정본 사건은 원인을 "재촬영 습관"이라 적지만 D+284의 회사는 모른다 — Day7 로그에서 안다. 예상 소요 6분.}
\end{frame}

\begin{frame}{ARR — 4,860과 6,060, 구축비 1,200만이 들어가면 분모가 바뀐다}
\fig[0.58]{fig_4_4_arr.pdf}
\keymsg{규칙을 적은 사람이 보고서에 6,060을 썼다 — 방어는 하나, 보고서에 두 숫자를 다 쓴다}
\note{그림 4-4 · 교재 4.4절 "작은 숫자로 먼저". 3사 × 3.0 × 45만 × 12 = 4,860(G-05) + 구축비 300/400/500 = 1,200(G-04) = 6,060(G-03 "계약 기준"). AI 원가 210 ÷ 405 = 51.9 / ÷ 505 = 41.6(G-09). GM 41(계약·예상 원가) / 26(순SaaS) / 12월 실제 청구서 반영 시 $-$72\% — 세 숫자가 전부 산수로 맞다, 정의가 다를 뿐. Day7에 열 배 큰 숫자로 돌아온다(F-03·F-04) — 오늘은 예고만. 예상 소요 8분.}
\end{frame}

\begin{frame}{계약의 다섯 칸 — 그리고 가격표의 제외 칸}
\begin{tbl}[\footnotesize]{L{24mm}L{62mm}L{40mm}}
조건 & 문구 & 근거·리스크 \\ \midrule
구축비 분리 & 별도 청구 · 서비스 매출 계정 · \textbf{ARR·MRR에 넣지 않는다}(1,200만) & bookings ≠ ARR. 리스크: "계약 ARR"이라는 말 \\
데이터 이용 & 비식별 검사 사진·기록만 · \textbf{도면 제외} · 미업로드 모드 & B-19 · 김도윤 공문 \\
판정 오류 책임 한계 & 연 계약액의 10\% = \textbf{162만} — ㉗ 항목 6과 동일 문장 & 클레임 2,700만과의 격차 — 고객은 서명하고 잊는다 \\
해지·환불 & 12개월 · 30일 통지 · 잔여 50\% · 구축비 환불 없음 & 12개월 미만 이탈 시 CAC 회수 불가 \\
가격 인상 & 연 1회 · 10\% 이내 · 90일 전 통지 & 공정 사용 조항의 도입 경로 \\
\end{tbl}
\keymsg{가격표의 제외 칸이 Won't 목록의 공개판이다 — 온프렘·MES·자동판정·포털 2종·항목 구조}
\note{교재 4.5절 · 워크북 ㉘ 항목 6·7. 책임 한계의 리스크 문장이 이 교재가 D+300에 할 수 있는 마지막 문장 — 그 격차가 무엇을 하는지는 Day7. Day8에 이 다섯 칸을 온담이라는 발주자에게서 받는다(배상 상한 10 vs 100\%). 4.6절(심화) 표는 언급만. 예상 소요 7분.}
\end{frame}

\begin{frame}{텀시트를 읽는 법 — 조항마다 "우리 숫자로 무슨 뜻인가", 없는 조항도 행으로}
\begin{tbl}[\footnotesize]{L{26mm}L{30mm}L{70mm}}
조건 & 값 & 우리에게 무슨 뜻인가 \\ \midrule
투자 / pre·post & 12억 RCPS / 48 · 60억 & 주당 4,800원. "회사 가치"가 아니라 이 라운드의 가격. D+149 잔액 1,505만 \\
청산우선권 & 1x 참가적, 상한 없음 & 12억 먼저 + 나머지의 20\%를 또. 100억에서 29.6(비참가적 20.0) \\
상환권 & 3년 후 연 3\% 복리 & 13.1억을 요구할 수 있으나 이익잉여금 한도 — 실질은 지렛대 \\
drag-along & FD 60\% & 노들 20 + 강민수 36 = \textbf{56\% < 60\%} — 창업자 둘 이상 필요. 창업 3인 72\%만으로 발동 가능 \\
TIPS & 8억(2년) · 유료 5개사(D+360) & R\&D 자금, 희석 없음 — 조건이 제품을 누른다 \\
\textbf{옵션풀 위치} & \textbf{조항 없음} & 협상 1번 · 구두 합의 · 문서에 없다 — 다음 리드가 정한다 \\
\end{tbl}
\keymsg{economics(돈이 어떻게 나뉘는가)와 control(누가 정하는가) — Feld \& Mendelson의 구분}
\note{교재 5.1절 · 워크북 ㉚ 항목 1. "뜻" 칸은 항목 3의 산수를 한 뒤에만 쓸 수 있다 — 항목 1을 마지막에 채운다. 서영채가 먼저 물을 것은 월간 보고의 ARR 정의. 예상 소요 8분.}
\end{frame}

\begin{frame}{캡테이블 T0 → T1 — 주식수로 계산한다}
\fig[0.56]{fig_5_1_captable.pdf}
\keymsg{48억 ÷ 1,000,000 = 4,800원 · 12억 ÷ 4,800 = 250,000주 · 36 / 24 / 12 / 8 / 20. 10 → 8\%는 감소가 아니라 희석}
\note{그림 5-1 · 교재 5.2절 "작은 숫자로 먼저". \%만 적은 캡테이블은 다음 라운드에서 계산할 수 없다. pre와 post의 주당가가 같은 것이 정상 — 옵션풀을 pre에 신설하면 달라진다(다음다음 프레임). 옵션풀 행의 비고 "추가 풀 없음, 위치 미정"이 다음 라운드로 넘어가는 문장. 예상 소요 6분.}
\end{frame}

\begin{frame}{"참가적" 한 단어의 값 — 100억에서 9.6억}
\fig[0.56]{fig_5_2_waterfall.pdf}
\keymsg{비참가적 = max(12, 20\% × 매각가) · 참가적 = 12 + 20\% × (매각가 $-$ 12). 작은 매각일수록 비중이 크다(30억에서 12\%)}
\note{그림 5-2 · 교재 5.3절. 30억: 12.0 vs 15.6 / 100억: 20.0 vs 29.6 / 300억: 60.0 vs 69.6. 창업자 3인(참가적·풀 pre 70\%) 12.6 / 61.6 / 201.6. 옵션풀 전량 행사 가정(보수적). 6교시 ㉚ 항목 3에서 학생이 100억 열을 네 조합 전부 직접 계산한다 — 그 산수를 한 사람만 협상 2번(참가 상한 3x 또는 비참가)을 올린다. 딥게이지는 올렸고 거절당했고 서명했다 — 왜는 BATNA. 예상 소요 8분.}
\end{frame}

\begin{frame}{옵션풀의 위치 — 금액이 아니라 선례가 값이다}
\fig[0.56]{fig_5_3_pool_position.pdf}
\keymsg{시드에서 0.43\%p, 시리즈A 15\% pre에서 1.29\%p — 딥게이지 계약에는 이 위치를 정하는 조항이 없다(D+186)}
\note{그림 5-3 · 교재 5.4절. 풀 10\% 요구 시 post 기준 투자자 19.57\%/창업자 70.43\%, pre 기준 20.0\%/70.0\%(주당 4,667원). 협상 1번 문구("라운드 후 전체 주주가 지분 비율대로 부담") → 서영채 "그건 시리즈A 때 얘기죠" → 구두 → 미기재 → D+186 발견 → 노들 "시리즈A 리드가 정할 사항". 부속 메모의 마지막 줄 — 구두 합의는 조항이 아니다. F-01의 두 번째 장면, 세 번째는 Day8 ㊴. 예상 소요 7분.}
\end{frame}

\begin{frame}{TIPS의 압력 · 협상 요청 3 · 약한 BATNA를 적는 정직함}
\begin{itemize}
\item \textbf{TIPS 조건은 새 결정을 요구하지 않는다} — 지킬 것인가를 묻는다. 라인 1개 계약의 유혹 → 최소 라인 2 유지 / 온프렘 고객의 유혹 → Won't 유지 / 지표 축소의 유혹 → G-18 병기
\item \textbf{협상 3} — 1 옵션풀 post-money 문구(← ㉑ 항목 3) / 2 참가 상한 3x 또는 비참가(9.6억) / 3 drag에 창업자 과반 동의. 결과: 구두 → 미기재 / 거절 / \textbf{수용}
\item 3번은 노들에게 비용이 없었다 — 그래서 수용됐고, 그래서 1번의 구두 합의가 생겼다. 상대에게 싼 것을 하나 섞는다
\item \textbf{BATNA — 약하다.} 다른 운영사 2곳 조건 미제시 · 잔액 1,505만 · TIPS 추천이 노들에 묶임. 9.6억은 100억 매각의 이야기, 1,505만은 다음 주 급여의 이야기
\end{itemize}
\keymsg{약한 BATNA를 적는 것은 패배 선언이 아니라 — 왜 2번 없이 서명했는지를 6개월 뒤의 자신에게 설명하는 기록이다}
\note{교재 5.5\til{}5.6절 · 워크북 ㉚ 항목 4·5. Fisher \& Ury BATNA. Day8에 시리즈A 텀시트 두 장을 받을 때 BATNA가 다르다 — 그때 이 문단을 다시 읽으라. 예상 소요 7분.}
\end{frame}

\begin{frame}{5교시 핵심 정리}
\begin{enumerate}
\item 청구서는 금액과 고객을 말하고 원인을 말하지 않는다 — 가설 넷, 판별 계측, 첫 조치는 계측. 소급 없음, 상한 없음
\item ARR에 구축비를 넣지 않는다 — 4,860 vs 6,060. 보고서에 두 숫자를 다 쓴다
\item 계약 다섯 칸 — 책임 한계 162만은 ㉗과 같은 문장. 가격표의 제외 칸이 Won't의 공개판
\item 텀시트 — 뜻 칸을 우리 숫자로, 없는 조항도 행으로. 캡테이블은 주식수로(4,800원 × 250,000)
\item 참가적 = 100억에서 9.6억 · 옵션풀 위치는 선례가 값 · 협상은 문장으로 셋, BATNA는 정직하게
\end{enumerate}
\keymsg{구두 합의는 조항이 아니다}
\note{교재 제4·5장 핵심 정리 발췌. 예상 소요 3분.}
\end{frame}

% =====================================================================
\section{6교시 (75분) — 실습 ②: ㉘ 프라이싱·UE v1 · ㉙ PR/FAQ·사전등록 · ㉚ 시드 텀시트·캡테이블 v1}
% =====================================================================
\begin{frame}{이 교시에서 하는 것}
\begin{itemize}
\item \textbf{㉘ 프라이싱·밸류메트릭·UE v1} — 항목 1(4안 비교) · 2(선택과 대가) · 4(원가 붕괴 — 가설·계측·대응) — \textbf{30분}
\item \textbf{㉙ PR/FAQ·사전등록·베타 기준} — 항목 3(내부 FAQ 중 5문, 어려운 것부터) · 4(사전등록 X-04·X-07) — \textbf{20분}
\item \textbf{㉚ 시드 텀시트·캡테이블 v1} — 항목 2(캡테이블 주식수) · 3(매각가 — 100억 열 네 조합 전부) · 5(협상 3 + BATNA) — \textbf{25분}
\item 세 문서가 서로를 인용한다 — ㉘ 항목 6 ← ㉗ 계약 문구, ㉙ 항목 5 ← ㉖ 게이트, ㉚ 항목 5-1 ← ㉑ 항목 3
\end{itemize}
\keymsg{㉘ 항목 4는 정답이 없다 — 여러분의 가설 목록에 "계측"이 있는지만 확인한다}
\note{30 + 20 + 25 = 75. 카드가 주는 것 — 파일럿 3사 제원, 청구서 내역, X-07 표본 177, 텀시트 세부, 매각가 3구간. 예상 소요 3분.}
\end{frame}

\begin{frame}{㉘ 프라이싱·밸류메트릭·UE v1 — 과제 지시}
\begin{itemize}
\item \textbf{과제}: 항목 1(좌석/라인/건수/성과과금 × 5기준 H/M/L + 한 줄 근거) · 2(선택과 \textbf{대가}) · 4(v1.1 — 가설 셋 이상, 판별 계측, 조항·상한·경고). 항목 3의 산수는 §3.5로 확인하되 두 열의 공헌이익률은 직접 계산
\item 시점 v1.0 \textbf{D+137} / v1.1 \textbf{D+284}. 카드 — 정본 §3.2 UE 표, 공장당 라인 3.2, G-18 계약 3.0/실사용 2.4, 청구서 서비스별(612/41/58/69)과 테넌트별 합계(동보 444), 파일럿 3사 제원(라인 3 · 구축비 300/400/500 · 동보 사진 19,400), 책임 한계 162만
\item 앞 산출물 — ㉕ 커널 행동 2, ㉔ A-06, 라인당 39.3만/월, ㉖ Won't(가격표 제외 칸), ㉗ 항목 6 문구
\item 여러분이 결정하는 것 — H/M/L과 근거 · 대가 문장 · 가설과 계측 · 공정 사용·경고 값
\end{itemize}
\keymsg{탈락 — 4안 비교에 "원가 정렬(우리)" 기준이 없다 · 대가 칸이 비어 있다 · 가설 하나를 사실로 적고 가격만 올렸다}
\note{워크북 §3.1\til{}3.3 + ◆ 카드 + §3.6 판정 규칙. 라인 대신 건수를 고른 조 — 예측 가능성·판매 용이성의 대가를 적었으면 정답 처리. 선택의 대가(점포 등 Day8)는 말하지 않는다. 예상 소요 5분.}
\end{frame}

\begin{frame}{작성 요령 ㉘ — 4안 비교 · 대가 · 가설 넷 (항목 1·2·4)}
\begin{itemize}
\item \textbf{항목 1}: 기준마다 누구 관점인가(예측 가능성은 고객, 원가 정렬은 우리). 종합은 합산이 아니라 "어느 기준을 포기했는가"의 문장
\item \textbf{항목 2}: 근거에 숫자(39.3만/월, 클레임 1건 = 20개월치). 대가는 항목 4로 이어지는 문장 — "사진 수 편차가 큰 고객에서 변동원가율 4배"
\item \textbf{항목 4}: 가설마다 "어느 계측이 있으면 판별되는가". 대응은 판별 뒤의 것과 판별 전에도 할 수 있는 것(조항·경고)을 나눈다. 기존 계약에 소급하지 않는다. 상한을 두면 FN 리스크를 적는다
\item 구축비는 "첫 해 구독료"가 아니다 — 항목 6에서 회계 규칙까지
\end{itemize}
\keymsg{항목 4에서 도달할 결론 — "계측이 없으면 어느 가설인지 알 수 없다"}
\note{워크북 §3.3 항목 1·2·4 가이드. 예상 소요 4분.}
\end{frame}

\begin{frame}{㉙ PR/FAQ·실험 사전등록·베타 기준 — 과제 지시}
\begin{itemize}
\item \textbf{과제}: 항목 3(내부 FAQ 15문 중 어려운 5문 — 답이 없으면 "미확정 — X-nn") · 4(사전등록서 — X-04 판정 초안 베타, X-07 조도 경고: OEC·최소 표본·근거·분석 시점·성공/중단)
\item 시점 \textbf{D+158}. 카드 — X-07 효과 크기 18\% → 8\%, 최소 표본 arm당 177(두 비율 검정) → 200; X-04 표본 설계 600건(검사원 3 × 시간대 2 × 100, 설계효과 3), 분석 D+290; 게이트 숫자는 ㉖ 항목 7과 동일
\item 앞 산출물 — ㉔ 카드 ①②③(사전등록 문법), ㉖ 항목 5·7·8, ㉗ 항목 6, ㉘ 항목 6·7
\item 여러분이 결정하는 것 — 어느 질문이 미확정인가 · OEC 문장 · 중단 기준 · 되돌림 조건
\end{itemize}
\keymsg{탈락 — 15문 전부에 답이 있다 · OEC가 셋 · 분석 시점이 "결과가 나오면" · ㉖ 항목 7과 다른 숫자}
\note{워크북 §4.1\til{}4.3 + ◆ 카드 + §4.6 판정 규칙. 미확정이 넷보다 많아도 된다. X-07 표본을 300으로 적은 조(근거 있으면) 정답 처리. 예상 소요 4분.}
\end{frame}

\begin{frame}{작성 요령 ㉙ — 미확정 · OEC 하나 · 중단 기준 (항목 3·4)}
\begin{itemize}
\item \textbf{항목 3}: 어려운 질문부터 — FN 8.1\% 미합격, 청구서(D+284 추가), 조현우 이후, 유료 5개사, 원가. 답이 ㉖·㉗·㉘의 칸에 있으면 그 칸을 인용(새로 쓰지 않는다)
\item \textbf{항목 4}: OEC는 하나 — 두 지표면 "그리고"로 묶어 하나의 판정. 불확정의 처리를 미리 적는다(X-04: 수용률만 충족 → 초안 유지, A-05b 열어 둠). 중단 기준은 가드레일(회피 20\% → 초안 OFF)
\item 최소 표본 — 계산할 수 있으면 근거(효과 크기·α·검정력), 없으면(모집단 3) "전수"와 3/3
\item 분석 시점은 날짜로 고정하고 "중간에 보지 않는다"를 적는다 · SRM 점검을 적는다
\end{itemize}
\keymsg{결과를 보기 전에 서명한다 — 성공 기준을 결과 뒤에 적으면 모든 실험이 성공한다}
\note{워크북 §4.3 항목 3·4 가이드 + §4.5 예시본 항목 4. 예상 소요 3분.}
\end{frame}

\begin{frame}{㉚ 시드 텀시트·캡테이블 v1 — 과제 지시}
\begin{itemize}
\item \textbf{과제}: 항목 2(캡테이블 T0 → T1 — 주식수로) · 3(매각가 시나리오 — 최소 100억 열은 네 조합 전부, 산식을 옆에) · 5(협상 요청 3 — 텀시트 문장으로 — + BATNA)
\item 시점 v1.0 \textbf{D+172} / 부속 메모 D+186. 카드 — 정본 §3.2·K-01·K-02(12억 · pre 48 · 참가적 1x · 상환권 · drag 60\%), 텀시트 세부(참가 상한 없음, 투자자 동의 조항 없음, 이사회 3 + 옵저버, \textbf{옵션풀 위치 조항 없음}), 매각가 30/100/300억, 가상 비교 전제(풀 10\%: post 19.57/70.43, pre 20.0/70.0), 잔액 1,505만, 대안 운영사 2곳
\item 앞 산출물 — ㉑ 항목 3(옵션풀 post 요구)·4(베스팅)·8(런웨이), PM ⑮ 조달의 "조건 / 대안 / 결렬 시"
\item 여러분이 결정하는 것 — 각 조항의 "뜻" · 협상 순위와 문구 · BATNA 평가 · 다음 라운드 전제
\end{itemize}
\keymsg{탈락 — \%만 있고 주식수가 없다 · 참가적을 "1x만"으로 · 요청이 "검토 요청" · BATNA가 "다른 투자자"뿐 · 옵션풀 위치가 어느 행에도 없다}
\note{워크북 §5.1\til{}5.3 + ◆ 카드 + §5.6 판정 규칙. 협상 순위가 예시본과 달라도(2번을 1번으로) 근거가 항목 3의 숫자면 정답 처리. 예상 소요 5분.}
\end{frame}

\begin{frame}{작성 요령 ㉚ — 주식수 · 네 조합 · 협상 문장 (항목 2·3·5)}
\begin{itemize}
\item \textbf{항목 2}: 주당가 = pre ÷ T0 FD = 4,800원 → 신주 250,000 → 지분은 그 뒤. 옵션풀 행 비고에 "추가 풀 없음 / 위치 미정"
\item \textbf{항목 3}: 비참가적 = max(12, 지분 × 매각가) / 참가적 = 12 + 지분 × (매각가 $-$ 12). 풀 pre/post는 투자자 지분(20.0 vs 19.6\%)을 바꾼다. 실제 K-02(추가 풀 없음)는 별도 행. 옵션풀은 전량 행사 가정
\item \textbf{항목 5}: 요청은 텀시트에 들어갈 문장으로. 상대에게 싼 것을 하나 섞는다(drag). 결과를 적는다(구두 → 미기재 / 거절 / 수용). BATNA는 정직하게 — 약하면 약하다고
\item 부속 메모(D+186) — "구두 합의는 조항이 아니다"를 항목 1의 마지막 행과 항목 5의 1번에 기록
\end{itemize}
\keymsg{100억 열을 손으로 계산한 사람만 "참가적" 한 단어를 협상 목록에 올린다}
\note{워크북 §5.3 항목 2·3·5 가이드 + §5.5 예시본. 예상 소요 4분.}
\end{frame}

\begin{frame}{예시본 참조 — 워크북 §3.5 · §4.5 · §5.5}
\begin{itemize}
\item \textbf{㉘} 라인 45만(원가 정렬 M을 알고 포기) · 두 열 86.6 / 48.5\% · v1.1 항목 4: 가설 H1\til{}H4 + 판별 계측 + 첫 조치 계측 · 공정 사용 6,000건·40원 · NDR 상한 105\% · 계약 다섯 칸 · 가격표 제외 칸
\item \textbf{㉙} 보도자료 8칸 · 고객 FAQ 12 · 내부 FAQ 15(미확정 4) · 사전등록 4(X-04 600건 D+290 · X-05·06 전수 · X-07 177 → 200 D+280) · 게이트 + 되돌림
\item \textbf{㉚} 여덟 조항 + 옵션풀 위치(없음) · 1,250,000주 · 4,800원 · 네 조합 × 3구간(100억: 20.0 / 29.6) · TIPS 압력 3 · 협상 3(구두 / 거절 / 수용) · BATNA 약함 · 부속 메모 D+186
\item 대조 규칙 — 칸별로. ㉘ 항목 4는 "계측"의 유무만 · ㉙ 항목 3은 미확정의 유무만 · ㉚ 항목 3은 100억 열의 산식만
\end{itemize}
\keymsg{인쇄용 완성 예시 — 00\_Day6\_템플릿팩\_완성예시.pdf}
\note{예시본은 저녁에. 수업 중에는 판정 규칙 상자만. 예상 소요 3분.}
\end{frame}

\begin{frame}{6교시 정리 — 다섯 문서가 한 사슬이다}
\begin{itemize}
\item ㉓ 트리 → \textbf{㉖} 백로그(유일한 상위 문서) → 항목 6 DoD → \textbf{㉗} 합격 기준 → 계약 문구 → \textbf{㉘} 항목 6(162만) → 구축비 분리 → (Day7 ㉝ 분모)
\item ㉖ 게이트 = \textbf{㉙} 항목 5 · ㉖ 리스크 5 = ㉙ 미확정 4 + 사전등록 4 → (Day7 ㉛ 판독 규율)
\item ㉑ 항목 3 → \textbf{㉚} 협상 1번 → 구두 → 미기재 → 부속 메모 → (Day8 ㊴ 1.29\%p)
\item 세 사건이 세 v1.1 — 퇴사(㉖) · 청구서(㉘) · 조항 부재(㉚). 사건 하나에 문서 전부를 고치지 않는다
\end{itemize}
\keymsg{범위 시점에 잘못 적힌 것은 측정 시점에 고쳐지지 않는다 — 그래서 오늘 다섯 문서에 그렇게 많은 칸이 있었다}
\note{워크북 §6 상호 관계 표. 예상 소요 3분.}
\end{frame}

% =====================================================================
\section{마무리 (30분) — 학술 배경 · 읽을 자료 · Day7 예고}
% =====================================================================
\begin{frame}[shrink=10]{학술 배경 — 깊이 읽기 (과정 후)}
\begin{tbl}[\footnotesize]{L{34mm}L{60mm}L{48mm}}
문헌 & 무엇인가 & 오늘의 어느 절 \\ \midrule
Singer, Shape Up (2019) & appetite · 회로 차단기 · scope hammering & 2.1\til{}2.4 — 왜 기간 고정인가 \\
Kohavi 외 (2020) · Nosek 외 (2018) & OEC · SRM · peeking · 사전등록의 논리 & 2.8 — 표본이 작아도 규율은 같다 \\
Shankar 외 (2024) · Gebru · Mitchell & criteria drift · 데이터셋·모델 문서화 & 3.3\til{}3.4 — 여집합, 개정 절차 \\
Beyer 외 SRE (2016) · ISO/IEC 42001 & error budget policy · 인간 감독 & 3.5\til{}3.6 \\
Skok · a16z · Sacks (2020) & ARR의 정의 · bookings ≠ ARR · burn multiple & 4.4 — 4,860 vs 6,060, 4.26 \\
Feld \& Mendelson (4판, 2019) · YC SAFE & economics vs control · 옵션풀 위치 & 5.1\til{}5.6 \\
\end{tbl}
\keymsg{문헌은 다툰다 — appetite(견적 없음)와 MoSCoW(견적 위의 비율)를 섞어 쓰는 이유를 알고 쓴다}
\note{교재 제6장 6.1\til{}6.5절 요약. 읽는 순서는 6.6절 표 — Shape Up → Feld → Bryar \& Carr → Skok·a16z → Kohavi·Nosek → SRE·Shankar → 표준. 판본은 인용 전에 확인(제작 노트 §1.2). 예상 소요 8분.}
\end{frame}

\begin{frame}{읽을 자료 — 필수 · 심화 · 참고}
\begin{itemize}
\item \textbf{필수(과정 전)} — Singer, \emph{Shape Up} Part 1\til{}2 · Bryar \& Carr, \emph{Working Backwards} 4장 · Feld \& Mendelson, \emph{Venture Deals} 4\til{}5장 · 정본 §3.2·§4.2\til{}4.4·§6·§7.3·§7.5 · Day5 워크북 ㉓ 항목 4·7, ㉕ 항목 6, ㉑ 항목 3
\item \textbf{심화(30일)} — Kohavi 외 1\til{}3·17장 · Skok "SaaS Metrics 2.0" · a16z "16 Startup Metrics" · SRE 3\til{}4장 / \textbf{(90일)} — Shankar 외 · Gebru 외 · Mitchell 외 · Wasserman 9\til{}10장 · Croll \& Yoskovitz · Moore 4장
\item \textbf{참고} — ISO/IEC 25010 · 42001 · Google Rules of ML · YC SAFE post-money · Sacks "The Burn Multiple" · PM 트랙 Day3 ⑫·⑮, Day4 ⑰, Day2 ⑥
\item \textbf{읽지 않아도 되는 것} — "MVP 만드는 법" 개론서(appetite와 견적의 구별이 없다) · 텀시트 표준 양식 모음(조항의 뜻이 없다) · 벤치마크 숫자만(분모의 정의가 없다)
\end{itemize}
\keymsg{워크북 X.5·X.6은 오늘 저녁, 필수 문헌은 Day7 전에}
\note{교재 6.7절. 예상 소요 5분.}
\end{frame}

\begin{frame}{Day7 예고 — PMF 추적: 곡선은 평탄해졌는가}
\begin{itemize}
\item D+300 → D+600, 데이터 일자 \textbf{2027-10-23}. 21명, 유료 31개사. 오늘의 "합격 전 베타"·"계측 먼저"·"6,060만"·"예산이 놓치는 FN"이 열 달 뒤 어떻게 되는지 — Day7 워크북 §0.3
\item ㉛ 코호트 리텐션 판독 — 삼각행렬 3장, \textbf{분모를 묻지 않으면 판독할 수 없다}. 오늘 X-04 사전등록의 규율이 그 판독의 규율
\item ㉜ North Star · ㉝ UE v2(같은 데이터의 두 답 — 오늘의 4,860 vs 6,060이 기준선) · ㉞ 사고 보고서·Eval v2·오류 예산 개정(오늘의 개정 절차가 시험받는다) · ㉟ 로드맵·예산 3분리(오늘의 Won't가 재검토된다)
\item 준비 — 오늘 쓴 ㉗(개정 절차·여집합)·㉘(항목 4·6)·㉙(사전등록)를 손 닿는 곳에
\end{itemize}
\keymsg{오늘 여러분이 적은 여집합의 첫 줄과 오류 예산이 놓치는 FN — 그 두 문장이 Day7의 사건이다}
\note{교재 "Day6를 마치며". D+600의 숫자는 오늘 말하지 않는다(제작 노트 §1.3). 예상 소요 5분.}
\end{frame}

\begin{frame}{Day6 핵심 정리 — 네 질문, 다섯 문서}
\begin{itemize}
\item \textbf{무엇을 만들지 않을 것인가} — Appetite 6개월 · 41 → 24 → 17 · 정의 축소 · 손실 분해 3.9 + 3.1 · Won't의 재검토 조건
\item \textbf{AI 기능은 언제 된 것인가} — 24/296과 173/1,184 · 모집단과 여집합 · FN 8.1\% 불합격, 기준 유지 · 예산이 놓치는 FN · 문구 둘
\item \textbf{무엇에 얼마를 받을 것인가} — 라인 45만, 원가 정렬 M을 알고 · 86.6 / 48.5\% · 780만의 가설 넷과 계측 · 4,860 vs 6,060
\item \textbf{누구 돈으로 어떤 조건으로} — 4,800원 × 250,000주 · 참가적 9.6억 · 없는 조항 · 약한 BATNA
\end{itemize}
\keymsg{범위 시점에 잘못 적힌 것은 측정 시점에 고쳐지지 않는다 — 오늘 여러분은 여집합을 비우지 않았고, 0.1\%p를 반올림하지 않았고, 구두 합의를 조항으로 세지 않았다}
\note{하루 마무리. 예상 소요 3분.}
\end{frame}
